\section{The critical annihilators for outer degrees four and five}
\label{sec:canonical-rank-laws}

\paragraph{R-NEW-002: a consequence of the established full rank laws.}
This section resolves the actual critical annihilator requested in
Question~6. It does not claim a new proof of the published fourth- and
fifth-degree base calculations. In particular, it does not identify the
entire larger common-polarization module $\mathcal N_r$ with the
annihilator that injects into it.

Use the normalized canonical grid map
\[
 h_{d,e}:\operatorname{Sym}^{d}(\operatorname{Sym}^{e}V)
       \longrightarrow\operatorname{Sym}^{e}(\operatorname{Sym}^{d}V),
 \qquad
 a_{d,e}^{\lambda}=[S_\lambda V:
             \operatorname{Sym}^{d}(\operatorname{Sym}^{e}V)].
\]
With the averaged symmetric-tensor evaluation pairings, the transpose
identity is exactly
\[
                  h_{d,e}(V)^*=h_{e,d}(V^*).
 \tag{R-NEW-2a}
\]
Indeed the row averaging idempotents are self-adjoint and the transpose
of the grid permutation is its inverse. This is the normalization
checked in GCT\_PLETH R-031, equations (OF.5)--(OF.8); no scalar is
suppressed.

For every partition $\lambda\vdash4n$, the image multiplicity is
\[
 [S_\lambda V:\operatorname{im}h_{n,4}]=
 \begin{cases}
 a_{n,4}^{\lambda},&1\le n\le4,\\
 a_{4,n}^{\lambda},&n\ge4.
 \end{cases}
 \tag{R-NEW-2b}
\]
The two formulas coincide at $n=4$. For every
$\lambda\vdash5n$, the corresponding complete law is
\[
 [S_\lambda V:\operatorname{im}h_{n,5}]=
 \begin{cases}
 a_{n,5}^{\lambda},&1\le n\le4,\\
 a_{5,5}^{\lambda}-\mathbf1_{\lambda\in\mathcal E},&n=5,\\
 a_{5,n}^{\lambda},&n\ge6,
 \end{cases}
 \tag{R-NEW-2c}
\]
where
\[
\begin{split}
\mathcal E=\{&(14,7,2,2),(13,7,2,2,1),(12,7,3,2,1),
 (12,6,3,2,2),\\
 &(12,5,4,3,1),(11,5,4,4,1),(10,8,4,2,1),(9,7,6,3)\}.
\end{split}
\]
All eight defects are one. The rank on the entire isotypic component
is this image multiplicity times $\dim S_\lambda V$. A shape of length
greater than $\dim V$ contributes zero. At $n=0$, the scalar map on
the empty partition is the identity.

\paragraph{Proof and exact primary inputs.}
Jacob (2004), Definition~4.1.1, p.~92, section~4.2, pp.~95--99,
and Table~4.4, pp.~104--106, computes the fourth-square transverse-grid
operator. Its collapsed adjacency-algebra matrix has determinant
$2^{223}3^{67}$ (p.~99), hence is invertible. The normalized
multilinear matrix of $h_{4,4}$ is $24^{-3}M_{4,4}$, as checked in
GCT\_PLETH R-031. Thus $h_{4,4}$ is invertible in every dimension.
M\"uller and Neunh\"offer (2005), section~4, final paragraph, explicitly
attribute this fourth-square calculation to Jacob.

McKay's propagation theorem, in Ikenmeyer's formulation (arXiv:
1509.04957, Theorem~1.3, p.~2; proof sections~2--6, pp.~3--8),
propagates injectivity of $h_{4,4}$ to $h_{4,n}$ for all $n\ge4$.
Transpose duality proves the second line of (R-NEW-2b).
For $n=2,3$, Sra (2026), Theorem~1.3, p.~2, gives surjectivity of
$h_{4,n}$ because $4\ge(n-1)^2$; transposition gives injectivity
of $h_{n,4}$. The $n=1$ map is the identity. This proves
(R-NEW-2b).

Cheung, Ikenmeyer, and Mkrtchyan (2015), Theorem~6(a),(b), p.~3,
prove injectivity of $h_{5,6}$ and the displayed multiplicity-free
kernel of $h_{5,5}$. McKay propagation and (R-NEW-2a) give
surjectivity of $h_{n,5}$ for every $n\ge6$. At $n=4$, injectivity
of $h_{4,5}$ is the fourth-degree result. At $n=2,3$, Sra's
same theorem applies to $h_{5,n}$, since $5\ge(n-1)^2$;
transpose again, and use the identity at $n=1$. This proves
(R-NEW-2c), the law already established in GCT\_PLETH R-031.

\paragraph{The requested critical submodules are exactly zero.}
Put $b=r^2-1$. The actual annihilator in GCT\_PLETH R-022,
equation (K.28), is
\[
 \mathcal K_{r+1,b}(W)=\ker h_{r+1,b}(W)
       \hookrightarrow\mathcal N_r(W)\otimes\operatorname{Sym}^{b}W.
\]
For $r=3$, $b=8$ and $h_{4,8}$ is injective by (R-NEW-2b)
and its proof. For $r=4$, $b=15$ and $h_{5,15}$ is injective
by (R-NEW-2c) and its proof. Therefore
\[
 \boxed{\mathcal K_{4,8}=0,\qquad\mathcal K_{5,15}=0.}
 \tag{R-NEW-2d}
\]
Their images in the two common-polarization tensor modules are the
zero submodules. In particular every isotypic component in the
previously isolated critical band has zero annihilator; by transposition,
the forward maps $h_{8,4}$ and $h_{15,5}$ have full target rank.
This conclusion uses the established full rank laws and is
stronger than merely excluding labels by the Pieri bounds. It does
not assert $\mathcal N_3=0$ or $\mathcal N_4=0$.

\paragraph{Comparison with the stated quadratic window.}
Sra's Theorem~1.3 concerns the actual map $h_{n,a}$ for
$n\ge(a-1)^2$. For $a=4$ its window starts at $n=9$;
for $a=5$ it starts at $n=16$. Every target isotypic component
is reached there. There is consequently no component whose rank
falls below the target multiplicity inside either window.
The eight defects at $(a,n)=(5,5)$ lie outside it. They refute
injectivity of that square canonical map, while the Foulkes
multiplicity inequality at equal parameters is equality.

\paragraph{Reverification and attribution.}
The full R-031 report and the actual injection (K.28) were read in
this run. The cited primary map definitions, the fourth-square
source and determinant, the fifth-square kernel, and the propagation
statement were checked. Sra's whole Lemma~2.1 and the complete proof
of Theorem~1.3, sections~2.1--2.2, pp.~3--5, were read and checked.
The retained exact Jacob matrix may be replayed using the supplied
GCT\_PLETH certificate. The large original Cheung--Ikenmeyer--
Mkrtchyan matrices are cited computations; they have not been
recomputed here. Formula (R-NEW-2d) is an immediate consequence
of these known results, not a historical novelty claim.

